\documentclass[10pt,a4paper]{amsart}
\usepackage[margin=19mm]{geometry}
\usepackage{amsmath,amssymb,mathtools}
\usepackage[scr=rsfs,cal=euler]{mathalfa}
\usepackage{xfrac}
\usepackage[unicode,hidelinks,bookmarks=false]{hyperref}
\newcommand{\ie}{i.e.\ }
\newcommand{\cf}{cf.\ }
\newcommand{\resp}{respectively\ }
\newcommand{\Subsection}{\subsection}
\providecommand{\nobreakdash}{\nobreak\hbox{-}}
\setlength{\emergencystretch}{2em}
\multlinegap0pt
\usepackage{bm}
\usepackage{bbm}
\usepackage{dsfont}
\usepackage{xspace}
\usepackage{cancel}
\usepackage{mathtools}
\usepackage{enumitem}
\usepackage{amsthm}
\makeatletter
\@ifundefined{theorem}{\newtheorem{theorem}{Theorem}}{}
\@ifundefined{lemma}{\newtheorem{lemma}[theorem]{Lemma}}{}
\makeatother
\newcommand{\ga}{\mathit{g}} % Safe math-italic for metric tensors
\title[The $G$-equivariant Kazdan--Warner problem]{The $G$-equivariant Kazdan--Warner problem}
\author{Leonardo F. Cavenaghi, Jo\~ao Marcos do \'O, Llohann D. Speran\c{c}a}
\date{Source: arXiv:2106.14709v5 (CC BY 4.0)}
\begin{document}
\maketitle

\noindent\emph{Source and licence.} The $G$-equivariant Kazdan--Warner problem. Version arXiv:2106.14709v5 (CC BY 4.0), distributed under the Creative Commons Attribution 4.0 International licence, as recorded in the arXiv licence metadata for this exact version and shown on the abstract page https://arxiv.org/abs/2106.14709v5. Full rights notice: licenses/arxiv-2106.14709-CC-BY-4.0.txt.\par\smallskip

\noindent\textit{Editorial scope.} Counted unit: the Lemma whose source label is \texttt{lem:dense\_kernel} in the article (source0.tex lines 656--658, source label \texttt{lem:dense\_kernel}), delivered together with the article's own complete original proof of it (source0.tex lines 659--710, one \texttt{proof} environment of 52 lines). No mathematics is written, repaired or re-derived: statement and proof are the article's own text line for line. Required context retained uncounted: eq:adjoint (lines 420--422); lem:marsden (lines 425--431). Every definition, display and label of those lines is retained unchanged. Omitted from this package and not redistributed: 1--419, 423--424, 432--655, 711--1279. Nothing inside a retained excerpt is omitted or altered: every retained body is delivered exactly as the article's lines are written, so no omission receipt of any kind accompanies this package. Citations: the retained excerpts carry no citation command at all, so no bibliography entry of the article is transcribed and no reference is redistributed. Reference adaptation: none was needed and none was made; every reference inside the delivered unit resolves inside it, so no unresolved reference or citation is produced. Printed slips kept literal: no printed word, symbol, hypothesis or number of the counted statement or of its proof was altered. Layout and class: the article's own class is \texttt{amsart} with packages including etex, fontenc, lmodern, babel, textcomp, amsmath, amsthm, amssymb, mathtools, mathrsfs, mathalfa, bbm, bm, stmaryrd; the wrapper is the lane's pinned \texttt{amsart} editorial wrapper at 10pt on A4, which restates the theorem-like environments the retained text needs and adds the article's own macro definitions that the retained text needs (\texttt{\textbackslash ga}), with extra wrapper packages (bm, bbm, dsfont, xspace, cancel, mathtools, enumitem). The delivered numbering is the unit's own. Assets: no retained span contains a figure environment, a graphics-inclusion command, a PDF-page inclusion, a table or a TikZ picture; no image file is copied into this package, the package declares no asset dependency and no image file is redistributed with it. Licence: version arXiv:2106.14709v5 of this article is distributed under the Creative Commons Attribution 4.0 International licence, as recorded in the arXiv licence metadata for this exact version and shown on the abstract page https://arxiv.org/abs/2106.14709v5. A full rights notice is delivered at licenses/arxiv-2106.14709-CC-BY-4.0.txt. Characters: every retained excerpt is pure ASCII, so the delivered engine, XeLaTeX with its default Latin Modern text font, reports no missing character.
\begin{equation}\label{eq:adjoint}
    A^*u = -(\Delta_{\ga} u)\ga + \nabla^2 u - u \mathrm{Ric}(\ga).
\end{equation}

\begin{lemma}\label{lem:marsden}
Let $\ga \in \mathcal{M}_G^{2,p}$ be smooth, and let $A^*$ be the formal adjoint \eqref{eq:adjoint} of the linearization of $F$ at $\ga$. Throughout, $\ker A^*$ refers to the $G$-invariant solutions of $A^* u = 0$.
    \begin{enumerate}[label=\textup{(\alph*)}, font=\upshape]
        \item If $\ker A^*$ is non-trivial, then $F(\ga)$ is a positive constant or $\mathrm{Ric}(\ga) = 0$. 
        \item If $\ker A^* = \{0\}$, then $L = AA^* : W_G^{4,p}(M) \to L_G^p(M)$ is an isomorphism. In particular, $A: W_G^{2,p}(M; \mathrm{S}^2T^{\ast}M) \to L_G^p(M)$ is surjective.
    \end{enumerate}
\end{lemma}

\begin{lemma}\label{lem:dense_kernel}
The set of $G$-invariant metrics $\ga \in \mathcal{M}_G^{2,p}$ for which $\ker A_{\ga}^* = \{0\}$ is open and dense in $\mathcal{M}_G^{2,p}$. In particular, if a smooth metric $\ga_0$ satisfies $C_1 < \mathrm{scal}_{\ga_0} < C_2$, there exists a smooth metric $\ga$ arbitrarily close to $\ga_0$ in $\mathcal{M}_G^{2,p}$ such that $\ker A_{\ga}^* = \{0\}$ and $C_1 < \mathrm{scal}_{\ga} < C_2$.
\end{lemma}

\begin{proof}
We have fixed everywhere $p>n$, so that metrics in $\mathcal M_G^{2,p}$ are of class $C^1$. For $\ga\in\mathcal M_G^{2,p}$ the operator
$A_\ga^*\colon W_G^{2,p}(M)\to L_G^p(M;S^2T^*M)$ is bounded and depends continuously on $\ga$ in operator norm. Its principal symbol
$\sigma_{A_\ga^*}(\xi)u=(|\xi|_\ga^2\,\ga-\xi\otimes\xi)u$ has trace
$(n-1)|\xi|_\ga^2\,u$, hence is injective for $\xi\neq0$. Thus, $A_\ga^*$ is
overdetermined elliptic and satisfies
$\|u\|_{W^{2,p}}\le C(\|A_\ga^*u\|_{L^p}+\|u\|_{L^p})$. Together with the compact
embedding $W^{2,p}\hookrightarrow L^p$, a standard argument gives
\[
  \ker A_\ga^*=\{0\}\iff \|A_\ga^*u\|_{L^p}\ge c\,\|u\|_{W^{2,p}}\ \text{for some }c>0 .
\]
Being bounded below is stable under small operator-norm perturbations and
$\ga\mapsto A_\ga^*$ is continuous, so $\{\ga:\ker A_\ga^*=\{0\}\}$ is open.

We claim that if $\ga'\in\mathcal M_G^{2,p}$ is smooth with
$\ker A_{\ga'}^*\neq\{0\}$, then there are smooth metrics $\ga_s'=(1+su)\ga'$, with $s$ arbitrarily small, such that $\ker A_{\ga_s'}^*=\{0\}$,
$\mathrm{scal}_{\ga_s'}\to\mathrm{scal}_{\ga'}$ uniformly, and
$\|\ga_s'-\ga'\|_{W^{2,p}}\to0$.

Indeed, by Lemma~\ref{lem:marsden}(a), we have that $\mathrm{scal}_{\ga'}\equiv\kappa\ge0$ is constant
($\kappa>0$, or $\mathrm{Ric}(\ga')\equiv0$ and $\kappa=0$). For $G$-invariant
$u$ and $|s|<\|u\|_\infty^{-1}$ the conformal metric $\ga_s'=(1+su)\ga'$ lies in $\mathcal M_G^{2,p}$, and, since $\left.\tfrac{d}{ds}\right|_{0}\ga_s'=u\ga'$,
the formula for $A$ gives
\[
  \left.\tfrac{d}{ds}\right|_{s=0}\mathrm{scal}_{\ga_s'}
  =A_{\ga'}(u\ga')=-\Delta_{\ga'}(n u)+\delta_{\ga'}\delta_{\ga'}(u\ga')-u\,\mathrm{scal}_{\ga'}
  =-(n-1)\Delta_{\ga'}u-\kappa u ,
\]
where we used $\delta_{\ga'}(u\ga')=-\mathrm{d} u$ and $\delta_{\ga'}(\mathrm{d} u)=-\Delta_{\ga'}u$.

Then, $\tfrac{d}{ds}|_{s=0}\mathrm{scal}_{\ga_s'}$ is constant only if $u\in E_0\oplus E_{\kappa/(n-1)}$  (for $\kappa=0$ only $E_0$), where
$E_\lambda\subset C_G^\infty(M)$ is the $G$-invariant $\lambda$-eigenspace of
$-\Delta_{\ga'}$. These eigenspaces are finite-dimensional. Since the action has cohomogeneity at least one,
the orbit space $X$ has positive dimension, so $C^\infty_G(M)\cong C^\infty(X)$ is
infinite-dimensional. Hence, we may choose $u$ outside this finite-dimensional subspace. Then $\left.\tfrac{d}{ds}\right|_{s=0}\mathrm{scal}_{\ga_s'}$ is
non-constant, hence $\mathrm{scal}_{\ga_s'}$ is non-constant for all small
$s\neq0$. Since Lemma~\ref{lem:marsden}(a) forces any metric with nontrivial kernel to have positive constant scalar curvature or to be Ricci-flat, both excluded, $\ker A_{\ga_s'}^*=\{0\}$. Finally, $\mathrm{scal}_{\ga_s'}\to\kappa$ uniformly and
$\|\ga_s'-\ga'\|_{W^{2,p}}=|s|\,\|u\ga'\|_{W^{2,p}}\to0$. This verifies the claim.


Given $\ga_0$ and $\epsilon>0$, pick a smooth $\ga_0'$ with
$\|\ga_0'-\ga_0\|_{W^{2,p}}<\epsilon/2$. If $\ker A_{\ga_0'}^*=\{0\}$ we are done.
Otherwise, the argument in the previous paragraph furnishes $\ga_s'$ with $\ker A_{\ga_s'}^*=\{0\}$ and, for
$|s|$ small and nonzero, $\|\ga_s'-\ga_0\|_{W^{2,p}}<\epsilon$.

 Let $\ga_0$ be smooth with
$C_1<\mathrm{scal}_{\ga_0}<C_2$. If $\ker A_{\ga_0}^*=\{0\}$, take $\ga=\ga_0$. Otherwise, we apply the claim above to $\ga_0$, obtaining smooth metrics $\ga_s=(1+su)\ga_0$, and set
$\epsilon_0=\min_{x\in M}\{\mathrm{scal}_{\ga_0}(x)-C_1,\,C_2-\mathrm{scal}_{\ga_0}(x)\}>0$.
For $|s|$ small and nonzero, $\ga_s$ is smooth, $\ker A_{\ga_s}^*=\{0\}$ and
$\|\mathrm{scal}_{\ga_s}-\mathrm{scal}_{\ga_0}\|_{L^\infty}<\epsilon_0$, whence
$C_1<\mathrm{scal}_{\ga_s}<C_2$. \qedhere
\end{proof}

\end{document}
